\scanpage{21}
% L01–04
des triples $(\alpha,R,\varphi_0)$, où a) $\alpha\in\Gamma(S,\mathcal O_S)$ tel que $\alpha+2$ inv (i.e. $\alpha(s)\ne-2\quad\forall s\in S$), b) $R$ un torseur sous \inserted{le groupe} $\deleted{\unsure{G}}\ U_{\beta=\alpha-2}$ c) $\varphi_0:\gamma\longrightarrow({}_2(U_\beta))(S)$ est un hom de groupes.

% L05–09
\uline{Proposition} Considérons \unsure{toujours|l'hom.|homomorphismes} \unsure{une|sur} \deleted{\unsure{$i$ entier pair,}} \deleted{\unsure{et soit}}

\deleted{\uline{Corollaire} Considérons $\Gamma=\mathbb D_i\times\mathbb Z_{/2}$ (\unsure{$i$ entier $\in(2\mathbb N)$}) et un corps alg. clos $k$, enfin la condition}

\[
\deleted{(\ast)\quad\exists\ \text{hom injectif}\quad\Gamma=\mathbb D_i\times\mathbb Z_{/2}\lhook\joinrel\longrightarrow\mathrm{GP}(1,k).}
\]
\deleted{a.) Si $i=0$, Alors $(\ast)$ est satisfaite sss}

% L10–15
\unsure{situation|solution} $\varphi:\mathbb D_\infty\longrightarrow G$ ($G$ forme de $\mathrm{GP}(1)$). \deleted{\unsure{Cas 1|Cond.}} Soit $n\in\mathbb N^{\deleted{\unsure{*}}*}$, \inserted{$n\geq3$}, considérons $\mathbb D_{n}\overset{\text{can}}{=}\mathbb D_i\simeq\mathbb D_\infty/\{\underline u^n\}$. Considérons les conditions (équivalentes)

a) $u^n=1$, et $\forall s\in S$, $u_s$ d'ordre \unsure{$n$} exact \deleted{[Si $n=0$, il faut dire : $u_s$ d'ordre $\infty$]}

b) $\varphi$ se factorise par un \uline{mono} $\mathbb D_{i_S}\lhook\joinrel\longrightarrow G$.

% L16–22
On veut les exprimer en termes de l'invariant $\alpha$. On trouve : \deleted{\unsure{on est dans les cas}}

\deleted{1°) Si $n=\infty$ : il f. et s. que $\forall s\in S$, on ait \unsure{que a) et c)} $\forall s\in S$ tel \unsure{1°) où}}

c) 1°) $n.1_S$ est inversible, \uline{ou} $n$ \inserted{est} \deleted{\unsure{$=2p$ où $p$|$=p$} \unsure{premier}} \inserted{un} \deleted{ou bien $n$ \unsure{d'un $p$ premier dans} $\Gamma(S,\mathcal O_S)$} est premier impair \deleted{\unsure{inversible}}, \deleted{\unsure{alr}} \uline{et} 2) $\Psi_n(\alpha)=0$

[$\Psi_n\in\mathbb Z[U]$ polynôme $\frac12$ cyclotomique d'indice $n$]
% archive sheet number: 20
